% TODO: fill in hardware numbers (Table I, and every XX% / XX.X placeholder
% below) once trial collection against the Legged Perceptive Stack completes.

\subsection{Hardware Comparison}
\label{subsec:results-hardware-comparison}

We evaluate GaitNet on a physical Unitree Go1 against the
\textit{Legged Perceptive Stack} \cite{leggedperceptive}, a
leading framework that is open source, combining NMPC and WBC atop
the OCS2 solver, using rigidly scheduled trajectories optimized in
phase. On hardware, GaitNet supplies footstep targets directly
to this same NMPC/WBC pipeline in place of its native gait scheduler
(\autoref{subsec:methodology-system-overview}), so the comparison
below holds the low-level tracking controller fixed and isolates the
contribution of the planning layer alone.

A physical test track (\autoref{fig:figure-test-environment}) mirrors
the terrain used in simulation. Narrow strips of rigid wooden ground
are interspersed with gaps of varying width, distributed
pseudorandomly. We characterize a track's difficulty as the
percentage of its area that \emph{cannot} be stepped on, ranging from
$\corlDifficultyMin\%$ (flat ground) to $\corlDifficultyMax\%$ (dense
gaps). Because both conditions share a controller, tracking error
would mostly reflect the shared low-level pipeline rather than the
planner. We instead measure overall mission autonomy directly: success
rate (completing the
traversal without the robot's base contacting the ground or drifting
off the track), total steps taken, and traversal time, each over
\corlHwTrialsPerCell{} independent physical trials per method and
difficulty level.

We compare four conditions, isolating the contribution of emergent
concurrency directly:
\begin{enumerate}
  \item \textbf{Legged Perceptive Stack:} the NMPC/WBC baseline
    described above \cite{leggedperceptive}.
  \item \textbf{GaitNet-Single:} GaitNet with the contact-count
    constraint tightened to permit only one leg in swing at a time,
    otherwise identical to the full model.
  \item \textbf{GaitNet-Trot:} GaitNet restricted to only swinging
    diagonally opposing leg pairs concurrently, matching a
    conventional trot's phase structure.
  \item \textbf{GaitNet (Full):} the proposed planner under its
    standard constraint of at most two concurrent swing legs, with no
    restriction on \emph{which} legs those are.
\end{enumerate}
A trial is successful only if the robot completes the full traversal
without falling or placing a foot in a removed terrain section.

\begin{figure}[t]
  \centering
  \includegraphics[width=\linewidth]{images/figures/test-environment.png}
  \caption{The Unitree Go1 platform navigating the physical testing
    track used for hardware evaluation. Terrain difficulty is the
    percentage of the track area that cannot be stepped on.}
  \label{fig:figure-test-environment}
\end{figure}

\begin{table}[t]
  \centering
  \caption{Hardware Metrics on the Unitree Go1, \corlHwTrialsPerCell{}
    Trials per Method}
  \small
  \begin{tabular}{lccc}
    \toprule
    \textbf{Method} & \textbf{Success} & \textbf{Time [s]} & \textbf{Steps} \\
    \midrule
    Legged Perceptive Stack \cite{leggedperceptive} & \corlHwLPSSuccess\% & \corlHwLPSTime{} & \corlHwLPSSteps{} \\
    GaitNet-Single & \corlHwSingleSuccess\% & \corlHwSingleTime{} & \corlHwSingleSteps{} \\
    GaitNet-Trot & \corlHwTrotSuccess\% & \corlHwTrotTime{} & \corlHwTrotSteps{} \\
    \textbf{GaitNet (Full)} & \textbf{\corlHwFullSuccess\%} & \textbf{\corlHwFullTime{}} & \textbf{\corlHwFullSteps{}} \\
    \bottomrule
  \end{tabular}
  \label{tab:hardware-metrics}
\end{table}

\autoref{tab:hardware-metrics} and \autoref{fig:data-hardware-survival}
summarize hardware results across the full difficulty range.
GaitNet (Full) achieves the highest success rate of the four methods,
with the gap widening substantially at high difficulty
($\corlHighDifficultyRangeLow$--$\corlDifficultyMax\%$).
The Legged Perceptive Stack's fixed phase schedule, and the
GaitNet-Trot and GaitNet-Single ablations, all increasingly commit to
footholds inside a scheduled or single-leg sequence that cannot be
revised once the terrain beneath it turns out to be invalid, whereas
GaitNet (Full) reevaluates every leg independently at every 40\,ms
tick and can commit a second leg mid-swing to recover from a
foothold the first leg is already unable to use productively. This
matches the ordering predicted by
\autoref{subsec:methodology-high-rate-selection}: GaitNet-Single, the
most constrained ablation, performs closest to the scheduled baseline,
GaitNet-Trot recovers some concurrency and some performance, and
GaitNet (Full) is uniformly best.

\begin{figure}[t]
  \centering
  \includegraphics[width=\linewidth]{images/data/experiments/Gaitnet
  - survival-curr - 69.4 - individual.png}
  \caption{Hardware success rate versus terrain difficulty for all
    four evaluated methods. \textit{Placeholder: simulation figure
    shown pending hardware data collection.}}
  \label{fig:data-hardware-survival}
\end{figure}

\subsection{Simulation Results (Preliminary Framework)}
\label{subsec:results-simulation}

Before developing the system described in
\autoref{sec:methodology}, our preliminary framework \cite{Sullivan2025}
evaluated cost-ablated GaitNet against an independent single-leg
baseline planner in NVIDIA Isaac Lab, across a grid of terrain
difficulties ($\corlDifficultyMin\%$ to $\corlDifficultyMax\%$, in
$\corlDifficultyStep\%$ increments) and commanded forward velocities
($\corlVelocityMin$ to $\corlVelocityMax$\,m/s), over
\corlSimTrialsPerCell{} trials of \corlSimEpisodeDuration\,s each per
cell. These results were collected under the discrete,
randomly sampled candidate distribution described in
\autoref{subsec:methodology-gaitnet-training}
(\corlTrainCandidatesPerLeg{} candidates per leg, cost-ablated) rather
than the continuous candidates refined by gradients used at
deployment in this paper
(\autoref{subsec:methodology-gaitnet-gradient-ascent}). We report
them here for context. We have not rerun this grid under the
refinement procedure introduced in this paper, so they should not be
read as validating the system described here.

In that configuration, GaitNet (Full) achieved a
$\corlResCostAblatedSurvivalRate\%$ mean survival rate, compared to
$\corlResBaselineSurvivalRate\%$ for the independent single-leg
baseline, a footstep selector using the lowest cost, operating one
leg at a time
with a fixed $\corlBaselineSwingDurationMs$\,ms swing duration,
structurally similar to the GaitNet-Single ablation introduced above
but without a learned value function. GaitNet maintained strong
performance for commanded velocities below
$\corlStrongPerfVelocityThreshold$\,m/s, while the single-leg
baseline's performance degraded once commanded velocity exceeded
$\corlVelocityMin$\,m/s and degraded further as terrain difficulty
increased. The performance gap narrowed only in the easiest condition
($\corlVelocityMin$\,m/s, flat ground), consistent with concurrent,
high-rate footstep selection becoming more valuable as task difficulty
increases.

\subsection{Emergent Concurrency Analysis (Preliminary Framework)}
\label{subsec:results-concurrency}

The central claim of this paper is that coordination between legs
emerges from high-rate selection of a single action rather than being
planned. Our preliminary framework \cite{Sullivan2025} first
demonstrated this mechanism in simulation, under the same discrete
candidate distribution described in
\autoref{subsec:results-simulation}. \autoref{fig:data-swing-schedule}
reproduces a representative swing schedule from that study. Below a
$\corlSwingScheduleTransitionMark$\,s mark, at a low commanded
velocity, the robot moves one foot at a time with no overlap between
swings. After the commanded velocity increases, the robot begins
committing a second leg while the first is still airborne, without a
strict alternating pattern of any kind. This transition is not
triggered by any explicit mode switch in GaitNet. The same value
function and the same per-tick decision rule govern both regimes, and
the increased concurrency emerges purely from more candidates crossing
the no-action threshold more often as commanded velocity rises.
Because this decision rule is unchanged by the continuous refinement
procedure introduced in this paper (\autoref{sec:methodology}), we
expect the same emergent behavior to hold at deployment. The hardware
comparison in \autoref{subsec:results-hardware-comparison} tests this
directly.

\begin{figure*}[t]
  \centering
  \includegraphics[width=0.85\textwidth]{images/data/swing-schedule.png}
  \caption{Example swing schedule. Each row represents a leg, with
    color indicating swing phase. Below
    \corlSwingScheduleTransitionMark\,s the robot is commanded at low
    velocity and steps one foot at a time; above
    \corlSwingScheduleTransitionMark\,s, at higher commanded velocity,
    swings overlap without a fixed alternating pattern.}
  \label{fig:data-swing-schedule}
\end{figure*}

In this same study, swing durations vary rather than following a
fixed schedule. \autoref{fig:data-swing-duration-histogram} shows a
peak between
$\corlSwingHistPeakMin$\,s and $\corlSwingHistPeakMax$\,s for typical
forward motion, with a secondary mode down to
$\corlSwingHistSecondaryMode$\,s used for situational, faster steps
such as recovering a marginal foothold. The mean swing duration
across training settles to $\corlResMeanSwingDuration$\,s with a
standard deviation of $\corlResStdSwingDuration$\,s, indicating the
duration head is in active use rather than collapsing to a constant.

\begin{figure}[t]
  \centering
  \includegraphics[width=\linewidth]{images/data/swing_duration_histogram.png}
  \caption{Histogram of swing durations selected by GaitNet during
  evaluation.}
  \label{fig:data-swing-duration-histogram}
\end{figure}

\subsection{Swing Duration Ablation (Preliminary Framework)}
\label{subsec:results-duration-ablation}

To assess whether this dynamic duration output is important, our
preliminary framework \cite{Sullivan2025} compared standard GaitNet
against a Duration-Ablated-GaitNet trained identically but constrained
to a fixed swing duration of $\corlResDurationAblationFixedDuration$\,s,
the approximate mean of the standard model's learned distribution. The
result was a null finding. The duration-ablated model achieved a
$\corlResDurationAblatedSurvivalRate\%$ survival rate, statistically
indistinguishable from the $\corlResGaitNetSurvivalRate\%$ achieved by
the corresponding standard GaitNet variant under matched conditions. We
do not read this as evidence that dynamic duration is unimportant in
general. More plausibly, the tested terrain and velocity ranges were
not demanding enough to expose a difference, and the reward function
does not explicitly incentivize using the full duration range beyond
what is needed to avoid termination. We report this honestly as a
limitation of that evaluation rather than removing the duration head
from the architecture used in this paper, since its cost is negligible
and it may prove useful at higher commanded velocities or on more
aggressive terrain than tested there.
